\begin{tresc}
Pani Ania kończy 1 marca przyszłego roku tyle lat, ile wynosi suma cyfr roku jej urodzenia. Jej młodszy brat Antek w tym samym dniu też kończy tyle lat ile wynosi suma cyfr roku jego urodzenia. O ile lat starsza jest pani Ania od swojego brata?
\end{tresc}
\begin{wariant}{A} 9 \end{wariant}
\begin{wariant}{B} 10 \end{wariant}
\begin{wariant}{C} 12 \end{wariant}
\begin{wariant}{D} 15 \end{wariant}
\begin{wariant}{E} 18 \end{wariant}
